\begin{annexitem}{annex:p5.direct-pa-axioms}{thm:p05-direct-pa-axioms}{Exact PA axiom recognition and presentation transport}
This proof uses the actual \(\mathsf{ND}_{\mathsf{FOL}}\) tree rules, including its
formula-specific labelled discharge and the all-variable freshness condition
for \(\exists E\).  All coding assertions in this item are statements
about natural numbers.  No arithmetic representing formula is being selected
or proved to satisfy D2 or D3 here.

\paragraph{The seven arithmetic sentences and literal conventions.}
Decode each \(q_j\) in the body using the constructor table.  Tags
\(0,1,2,3,4,5,7,10,11,12\) mean respectively variable, zero,
successor, addition, multiplication, equality, negation, implication,
universal and existential quantification.  The decoded formulas are
\[
\begin{aligned}
 &\forall x_0\,\neg(Sx_0=0),\\
 &\forall x_0\,\forall x_1\,(Sx_0=Sx_1\to x_0=x_1),\\
 &\forall x_0\,(\neg(x_0=0)\to\exists x_1\,(x_0=Sx_1)),\\
 &\forall x_0\,(x_0+0=x_0),\\
 &\forall x_0\,\forall x_1\,(x_0+Sx_1=S(x_0+x_1)),\\
 &\forall x_0\,(x_0\times0=0),\\
 &\forall x_0\,\forall x_1\,
                  (x_0\times Sx_1=(x_0\times x_1)+x_0).
\end{aligned}
\]
These are precisely the seven active \(\mathsf Q\) axioms under the
declared choice \(x=x_0,y=x_1\) for the source's schematic variable
names, with the displayed universal-closing order.  Their codes are the
fully specified constructor expressions in the body, not metavariables
standing for unspecified axiom codes.  Every one is a sentence.  In the
third axiom the existential variable is distinct from the outer variable;
the encoding does not accidentally bind the argument of the predecessor
assertion.

For the three two-variable arithmetic axioms, either order of their two
universal closures has exactly the same variable-occurrence set.  The
closure-transport proof below applies to such an order change as well.
Bound names in a displayed closed axiom may be renamed by the actual alpha
rule after making them fresh from its matrix.  For example, rename all
binders first to pairwise distinct variables absent from both expressions,
and then compare the binding positions; this yields an alpha derivation
between two closed renamings of the same binding pattern.  Such a derived
variant is not silently made one of the seven literal accepted constants.
The source-name convention is fixed before assigning the seven codes.

\paragraph{Substitution identities used by induction.}
For every formula \(\theta\), the specified substitution satisfies
\(\theta[x_i/x_i]=\theta\) literally.  Prove this by induction on
terms and formulas.  It is immediate at variables and commutes with the
nonbinding constructors.  At a binder for \(x_i\), substitution stops;
at another binder \(x_j\), \(j\ne i\) and
\(x_j\notin\operatorname{Var}(x_i)\), so it commutes without renaming.
These are all quantifier cases of the actual definition.

For the replacements \(0\) and \(Sx_i\) in the induction constructor,
no binder is renamed either.  The zero term has no variables.  The only
variable of \(Sx_i\) is \(x_i\): a binder for it stops substitution,
and every other binder is absent from that replacement term.  Thus the
third, least-fresh-variable clause is never invoked in either of these
two substitutions.  Structural induction gives
\[
\begin{aligned}
 \FV(\theta[0/x_i])&=\FV(\theta)\setminus\{x_i\},\\
 \FV(\theta[Sx_i/x_i])&=\FV(\theta),
\end{aligned}
\]
including when \(x_i\) is absent freely or occurs only under another
\(x_i\) binder.  In particular, substitution never changes a bound
occurrence that the induction schema leaves untouched.

Let \(\theta\) be coded by \(m\), and write \(\mathcal I(\theta,i)\)
for the formula coded by \(J(m,i)\).  Direct decoding gives
\[
 \mathcal I(\theta,i)=
 \bigl(\theta[0/x_i]\land
       \forall x_i(\theta\to\theta[Sx_i/x_i])\bigr)
         \to\forall x_i\theta.
\]
The preceding identities prove
\[
 \FV(\mathcal I(\theta,i))=\FV(\theta)\setminus\{x_i\},
 \qquad
 \operatorname{Var}(\mathcal I(\theta,i))
                       =\operatorname{Var}(\theta)\cup\{x_i\}.
\]
For the second equality, the unchanged occurrences of \(\theta\)
in the step antecedent and conclusion retain every old variable, the
quantifiers add \(x_i\), and the substitution clauses add no other
variable.  Canonical closure adds only free variables already in this
set, so it leaves that all-occurrence set unchanged.  This last equality
will be needed when replacing axiom leaves under \(\exists E\).

\paragraph{Primitive recursion and exact recognition.}
In a valid term or formula code, each variable index is below its variable
constructor code, or below the argument-list code of the binder at which
it occurs, and hence below the complete expression code.  This proves
the bound used by the scan defining \(F(m)\).  Membership in the already
computed finite free-variable list is primitive recursive.  Scanning
\(j<m\) in order includes each free index once, so it gives exactly the
increasing list claimed.  The invalid-input default is specified separately.

The closure operation \(U\) uses one reversal and exactly the input list's
length many constructor updates.  The two substitution operations used by
\(J\) are the previously proved primitive recursive operations.  Filtering
out \(i\), applying the fixed constructor expression for \(J\), and
wrapping the result by \(U\) are therefore primitive recursive.  This
proves totality and primitive recursiveness of \(\mathsf{IndAx}\) for
all numerical inputs, including its zero value on a nonformula matrix.

Every \(\mathsf{IndAx}(m,i)\) for a valid \(m\) is a sentence, by
the free-variable calculation.  Its leading quantifiers are exactly the
canonical parameter prefix; the next constructor is an implication.
Its right child is \(N_{11}(i,m)\).  Thus stripping the maximal leading
universal prefix recovers this implication, its right child recovers
\(i,m\), and the recognizer reconstructs precisely the original code.
It accepts every such induction axiom, including those with no parameters
and those with an absent induction variable.  The explicit fixed-constant
test accepts all seven arithmetic axioms.

Conversely, an accepted input is either one of those constants, or the
last test establishes literal equality with
\(\mathsf{IndAx}(m,i)\) for the recovered formula code \(m\).
There is no inference from semantic or alpha equivalence in this direction.
The recovered \(m\) is a proper subtree of the input, and \(i\) is a
proper constructor argument, so both are below the input code \(a\).
The stripping loop has at most \(a\) updates; the free-index scan has
bound \(m<a\).  The substitution traversals have the fixed primitive
recursive bounds already supplied in the syntax construction.  All other
tests are exact tags, arities and equalities.  Hence this is a primitive
recursive recognizer, with no unbounded search for a matrix or an induction
variable.  Numerical values produced during reconstruction may be large;
primitive recursion imposes no separate bound on their magnitude.

The two explicit left-side comparisons test the exact base and step
instances.  They reject a step using another binder, a changed matrix
copy, or a substitution that entered a shadowing binder.  Literal equality
of the prefix with \(F(m)\setminus[i]\) rejects a reordered, repeated,
omitted or extra parameter.  In particular a true or derivable noncanonical
closure is not automatically an axiom.  This proves the recognition and
sentence assertions of \cref{thm:p05-direct-pa-axioms}.

\paragraph{Actual ND transport between closing conventions.}
Let \(\eta\) be any formula, and let \(P,Q\) be two lists of variable
indices such that \(U(P,\ulcorner\eta\urcorner)\) and
\(U(Q,\ulcorner\eta\urcorner)\) code sentences.  From an assumption
of the first sentence, apply \(\forall E\) to its outside binder using
that very variable as the supplied term.  The literal identity
\(\theta[x_j/x_j]=\theta\), proved above, removes that binder without
changing the remaining matrix.  Repeat through \(P\) to obtain \(\eta\).
Now apply \(\forall I\) in the reverse order of \(Q\), obtaining its
closure with the prescribed outside-in order.  The only open assumption
is the original closed sentence, so each \(\forall I\) side condition
is satisfied.  This uses precisely \(\ell(P)+\ell(Q)\) quantifier
inferences and one assumption leaf, including either empty list.
Exchanging \(P,Q\) gives the reverse derivation.  Repeated or vacuous
binders, if present in a source closure, are removed and introduced by
exactly the same valid rules; they are not inserted into the canonical
axiom recognizer.

Apply this with \(\eta=\mathcal I(\theta,i)\).  The source induction
axiom is the closure of exactly this formula for its given \(\theta\),
induction variable and parameter list.  The canonical list is
\(F(\ulcorner\theta\urcorner)\setminus[i]\).  Thus each source
induction closure is derivable from the corresponding canonical axiom,
and the canonical axiom is itself an allowed source induction instance.
If a source fixes another permutation of exactly the free parameters,
the reverse quantifier derivation instead produces the canonical closure
from that fixed source representative.  The seven arithmetic axioms coincide
with the fixed representatives declared in the body; changing only their
two-variable closing order uses the same transport.

To transport a whole derivation, replace only its surviving open axiom
occurrences by these finite derivations of the same formulas.  All new
axiom leaves may be unmarked, and the inserted derivations have no discharge
rules.  Every old conclusion formula and rule parameter is retained.
For source-to-canonical induction transport, the new canonical leaf has
variable-occurrence set
\(\operatorname{Var}(\theta)\cup\{x_i\}\), which is a subset of that
of the source closure it replaces.  In the reverse direction choose the
source closure of exactly the free parameter set; then the occurrence sets
are equal.  The arithmetic-axiom closing-order changes have equal occurrence
sets as well.  Every new axiom leaf is a sentence.

These facts check every possible effect on an old inference.  Propositional,
equality, substitution and alpha rules retain their exact premise and
conclusion formulas.  An original surviving occurrence was not discharged
by any rule on its path to the root, and new leaves are unmarked, so no
new discharge occurs and no old nonvacuity witness is removed.  The
\(\forall I\) test still holds because the new open leaves are closed.
For \(\exists E\), the eigenvariable was absent from every old remaining
open formula; it remains absent from the all-variable subset at each
replacement leaf.  Its major formula and conclusion are unchanged.
Hence its stronger all-occurrence freshness test also survives.  The
quantifier rules inserted in the transport were verified in the preceding
paragraph.  This gives actual ND derivation transports in both directions,
not just a comparison of models or of individual axiom truth values.

The transport is effective and uses bounded operations on the input
derivation: parse each open axiom, recover its matrix and parameter prefix,
and insert the displayed finite quantifier chain.  Induction on the original
tree constructs the resulting tree.  If numerical records are desired,
list its nodes in postorder and recompute the premise indices.  Closure
chains have lengths bounded by the original formula-code and the recovered
free-index scan bounds.  The original finite row list bounds the number of
insertions.  More explicitly, for an input proof-record code \(p\), every
row formula code is below \(p\), and a replacement has at most \(2p\)
quantifier steps.  With \(n\) original rows, the resulting tree has at
most \(n(1+2p)\) rows.  Compute transformed child blocks in the original
postorder and concatenate them with shifted local premise indices.  Append
a copy of their parent row whose premise list names those rebuilt child
roots; its other fields are unchanged.  At an open axiom leaf use its
explicit quantifier chain instead.  This is an \(n\)-step accumulated-row computation with
the already defined primitive recursive list operations.  Malformed
records or unrecognized source axiom shapes may return zero.  Thus a
total primitive recursive implementation agrees with the transport on
every valid input; no unbounded search for an ND proof is involved.

This completes the derivability assertion of
\cref{thm:p05-direct-pa-axioms}.  The proof is for the literal seven
representatives fixed before coding and for the active induction scheme.
It does not equate differently serialized closed formulas as numbers.
\end{annexitem}

\begin{annexitem}{annex:p5.standard-pa-checker}{thm:p05-standard-pa-numerical-proof}{Numerical PA proof checking and its source bridge}
Use the recognition and transport just proved in
\cref{annex:p5.direct-pa-axioms}.
\paragraph{The direct proof checker is primitive recursive.}
All outer tag, list-length and equality tests are primitive recursive.
For \(n\) rows, the inherited procedure computes its row table by
primitive recursion on \(j<n\).  The fixed rule-arity list, in tag order,
is
\[
 0,2,1,1,1,1,3,1,2,1,2,1,1,1,1,1,2,1,0,2.
\]
The parameter-list lengths, in the same order, are
\[
 0,0,0,0,0,0,2,1,0,1,0,0,1,3,3,1,4,0,1,4.
\]
These lengths are supplemented by exactly the type, formula and side-condition
tests in the existing twenty-rule definition.  In particular, they are not
used alone as a substitute for those tests.  The parent counts are finite
sums over the premise lists, each of arity at most three.  The open lists
are computed by the existing concatenation and discharge operations.

Once the tree checks pass, the final open list has at most \(n\) entries,
all distinct assumption-leaf positions.  Apply the primitive recursive axiom
recognizer to their formula codes and multiply its zero-one answers.
An empty product has value one.  The remaining root and empty-field tests
are primitive recursive.  On malformed input an error flag returns zero;
each traversal still has its stated list or syntactic bound.  This proves
primitive recursiveness of the complete \(\mathsf{StdProof}_{\PA}\)
relation.  All rows are actual annotated trees, so no rule matrix or term
witness is inferred by searching over arbitrary formulas.

\paragraph{Both standard-model proof-meaning directions.}
Suppose \(\mathsf{StdProof}_{\PA}(p,\ulcorner\sigma\urcorner)\) holds.
The inherited row induction reconstructs exactly one finite ND tree with
root \(\sigma\), and its computed open lists are exactly its surviving
assumption occurrences.  Every one passes the direct axiom recognizer,
so the proved recognition theorem puts it in \(\PA^{\rm d}\).
This is a \(\PA^{\rm d}\)-derivation.  The preceding transport supplies
the required source \(\PA\)-derivation of \(\sigma\).

Conversely, take a source \(\PA\)-derivation of the sentence \(\sigma\).
Transport its surviving axiom leaves to the canonical presentation as
proved above.  Encode the resulting actual tree in postorder, retaining
its inference witnesses, variable indices, formula/term parameters and
discharge labels.  Every local check and parent count succeeds by the
same row induction, and the final open list contains only canonical axiom
leaves.  Put the record in \(N_{19}(N_{15}(R),0)\).  Its direct axiom
checks and root check succeed, giving the required standard proof code.
When the tree has no open assumptions, the axiom product is empty and the
same construction works.  A one-row \(=I\) proof of \(0=0\) is a
concrete instance of this case.  Extra certificate data are rejected even
then, because the second outer field must be zero.

\paragraph{Bridge to the c.e.-certificate convention.}
Fix an index whose domain is exactly the codes of the chosen source
\(\PA\) axiom set, as required for the earlier certified relation.
A successful certified code supplies an ND derivation with source axiom
leaves by the established Comp theorem and numerical proof equivalence.
Apply the just-described leaf transport and use outer tag \(19\) and
empty certificate field to obtain a successful direct code.  Where its
open axioms already use canonical representatives, only the outer tag and
certificate field need change; in general the closure chains must be
inserted as well.

Conversely, a successful direct code supplies the source derivation by
the other transport.  Each of its finitely many open source axiom leaves
lies in the domain of the fixed index.  The accepted bounded-verification
theorem gives an accepting history for each such leaf.  In their computed
open-list order, place these histories in \(W\) and use outer tag \(16\).
This is a successful earlier certified proof code.  With no open leaves,
take \(W=0\).  No bound on the search for those accepting histories is
claimed for an arbitrary chosen index; finite existence is all this direction
uses.  A compiler of certificates with PA-provable behavior is a separate
internal task.

Outer tags \(19\) and \(16\) are distinct under the injective constructor
coding, so the two positive proof-code sorts are disjoint.  Their existential
proof meanings agree as proved; their numerical predicates are not identical,
and this bridge does not assert arithmetic equivalence between any selected
representing formulas.

\paragraph{What remains for a standard arithmetic formula.}
The direct axiom recognizer fixes all numerical choices: matrix recovery,
substitution, parameter ordering, constant axioms and rejected shapes.
The direct proof procedure fixes the outer tag, empty certificate field,
all twenty row tests and final leaf test.  To use the graph-formula theorem,
one must now fix its finite description and the descriptions implementing
these library operations, and verify the relevant invariants inside \(\PA\).
The present proof establishes no PA theorem merely because one of these
numerical operations is primitive recursive.  The standard-formula selection,
internal modus ponens, and evidence-to-proof construction remain at that
subsequent scope.
\end{annexitem}
